% !TEX root = ../main.tex

%----------------------------------------------------------------------------------------
% CHAPTER 9 - IMPLEMENTACIÓ I PROVES
%----------------------------------------------------------------------------------------

\chapter{Implementació i proves}
\label{Ch:Implementation}

%----------------------------------------------------------------------------------------

\section{Implementació del pipeline ETL}

El pipeline ETL s'ha implementat com un conjunt de mòduls Node.js estructurats segons el patró de pipeline modular:

\subsection{Extracció}
Dos extractors independents (\texttt{EmbassamentExtractor} i \texttt{PrecipitationExtractor}) utilitzen \texttt{axios} per obtenir dades de les APIs Socrata de la Generalitat \cite{gencat_dades, socrata}. Suporten mode incremental (diari) i sincronització completa (\texttt{--full-sync}). La configuració es gestiona amb \texttt{dotenv}.

\subsection{Transformació}
Els transformadors normalitzen dates a ISO 8601, eliminen duplicats, calculen percentatges d'ocupació i classifiquen alertes de sequera (crític $<$20\%, avís 20-50\%, normal 50-75\%, òptim $>$75\%). Per a precipitació, s'implementa una finestra lliscant de 2 anys.

\subsection{Càrrega}
\texttt{DataLoader} desa les dades en fitxers JSON amb metadades, fusió incremental per evitar duplicats, i neteja automàtica de còpies antigues (més de 30 dies).

%----------------------------------------------------------------------------------------

\section{Implementació del frontend}

\subsection{Estructura}
Aplicació SPA amb React 18 \cite{react}, React Router 6 i gestió d'estat amb Zustand \cite{zustand}. Els components es carreguen asíncronament amb \texttt{React.lazy()} i \texttt{Suspense} per a code splitting.

\subsection{Gestió de dades}
\texttt{waterDataService.js} encapsula l'accés a dades mitjançant \texttt{axios}, proporcionant filtres per data i estació, agrupació per nom, i càlcul de colors d'estat.

\subsection{Visualitzacions}
\begin{itemize}
    \item \textbf{Recharts 2} \cite{recharts}: Gràfics de línies, dispersió i KPIs amb integració nativa amb React.
    \item \textbf{Leaflet 1} \cite{leaflet}: Mapa interactiu amb markers i popups informatius.
\end{itemize}

L'estil utilitza Sass \cite{sass} amb metodologia BEM i CSS Grid/Flexbox per a disseny responsiu mobile-first.

%----------------------------------------------------------------------------------------

\section{Automatització amb GitHub Actions}

Dos workflows coordinats \cite{github_actions}:

\textbf{ETL Pipeline (\texttt{etl-pipeline.yml})}: S'executa diàriament a les 02:00 UTC (cron) o manualment. Fa checkout, instal·la dependències, executa l'ETL, copia fitxers JSON a \texttt{frontend/public/data/}, i fa commit/push amb \texttt{[skip ci]} per evitar bucles.

\textbf{Desplegament (\texttt{deploy-frontend-pages.yml})}: S'activa a push a \texttt{main}. Construeix el frontend amb Vite \cite{vite} i desplega a GitHub Pages \cite{github_pages}.

%----------------------------------------------------------------------------------------

\section{Proves realitzades}

\textbf{Proves ETL}: Validació de cada fase (extractors, transformadors, carregadors) amb Jest, verificant estructures de dades, normalització, fusió i neteja.

\textbf{Proves frontend}: Verificació de \texttt{waterDataService}, components de visualització (Recharts, Leaflet) i navegació React Router.

\textbf{Proves d'integració}: Execució completa amb dades reals, simulació d'errors d'API i validació del flux complet.

%----------------------------------------------------------------------------------------

\section{Codi font i vídeo demostratiu}

\textbf{Repositori:} \url{https://github.com/AlexRuizMaso/opendata-water-visualization}

\textbf{Aplicació:} \url{https://alexruizmaso.github.io/opendata-water-visualization}

\textbf{Vídeo del codi (3' màx):} Inclòs al vídeo demostratiu de l'apartat 9.3.

%----------------------------------------------------------------------------------------